\section{Process Supervision：逐步验证推理}

\subsection{Outcome 与 Process 的差别}

假设一条推理包含步骤 $s_{1:T}$。Outcome supervision 只给完整轨迹标签 $R\in\{0,1\}$；process supervision 为每一步提供 $r_t\in\{0,1\}$ 或连续分数。

\videofigure{figures/outcome-supervision.jpg}{Outcome supervision 只在整条推理结束后提供标签。}{00:22:20--00:24:00}

\videofigure{figures/process-supervision.jpg}{Process supervision 精确定位错误步骤并改善信用分配。}{00:24:00--00:26:02}

\begin{importantbox}{Process supervision 的两个价值}
第一，它能在错误刚出现时剪枝，节省后续 token；第二，它减少“错误过程碰巧得到正确答案”的假阳性。对训练而言，步骤级 reward 也比单一终局 reward 提供更密集的学习信号。
\end{importantbox}

\subsection{PRM800K：人工步骤标签}

Let's Verify Step by Step 收集约 80 万个人工步骤标签。标注者阅读模型解答，在每个步骤后判断目前过程是否仍然正确。为了提高数据效率，方法使用 active learning，把最有区分度或最可能暴露模型弱点的样本优先送给人类。

\videofigure{figures/prm-labeling.jpg}{PRM800K 对模型生成的数学推理进行步骤级人工标注。}{00:26:02--00:29:08}

训练后，PRM 输出每一步的正确概率 $q_t$。整条轨迹可用最弱步骤、乘积或对数和聚合，例如：

$$
S(y)=\sum_{t=1}^{T}\log q_t.
$$

\begin{itemize}
    \item $q_t=P(s_t\text{ correct}\mid x,s_{1:t})$；
    \item $S(y)$：轨迹总体过程分数；
    \item 对数和会惩罚任一低概率错误步骤，同时避免直接连乘下溢。
\end{itemize}

\subsection{PRM 相比 ORM 的收益}

\videofigure{figures/prm-vs-orm.jpg}{PRM 在候选选择上优于 ORM 与 majority voting。}{00:29:08--00:30:44}

PRM 的收益在长链推理中特别明显，因为轨迹越长，局部错误被最终答案掩盖的机会越多。它还支持 beam search：每生成一段就评分，只扩展高分前缀。

\subsection{Active Learning 与大规模标注}

随机抽样会浪费大量人力在明显正确或明显错误的步骤上。Active learning 选择 verifier 当前最不确定、模型间分歧最大或最接近决策边界的样本。

\videofigure{figures/active-learning.jpg}{Active learning 提高步骤标注的单位人力收益。}{00:30:44--00:33:02}

\subsection{分布外泛化}

\videofigure{figures/prm-ood.jpg}{训练后的 PRM 在分布外数学数据上仍保持相对优势。}{00:33:02--00:35:20}

这种泛化说明 PRM 学到了一部分通用的数学步骤合理性，但不能据此认为它掌握了严格证明。自然语言推理中的隐含假设、跳步和符号歧义仍可能骗过 verifier。

\begin{warningbox}{Process reward 也会产生新的 reward hacking}
模型可能学会写出“每一步都很像教科书”的冗长轨迹以取悦 PRM，而不是更有效地解决问题。若 verifier 偏好格式、长度或常见措辞，RL 会放大这些偏好。步骤级监督提高了分辨率，但没有消除代理目标风险。
\end{warningbox}

\subsection{本章小结}

人工 process supervision 提供精确、密集且可用于搜索的验证信号，代价是昂贵的专家标注。下一步自然问题是：能否让模型和环境自动产生步骤标签？

